\documentclass{article}
\usepackage[UTF8]{ctex}
\usepackage{amsmath,amssymb,amsfonts}
\usepackage{enumitem}
\usepackage{geometry}
\geometry{a4paper, margin=1in}
% === 错题本排版优化（不参与题目标识） ===
\setlength{\parindent}{0pt}
\setlength{\parskip}{0.35em plus 0.10em minus 0.08em}
\setlength{\emergencystretch}{3em}
\setlist{itemsep=0.25em, topsep=0.40em, parsep=0pt}
\allowdisplaybreaks[2]
% === 错题本排版优化结束 ===
\begin{document}
\section*{暨南大学2026年高等代数真题}
\section*{149．暨南大学2026年高等代数试卷（回忆版）}
（注．题目来自考生回忆，部分表述或数据可能有误）

\subsection{题目1}
1．（10 分）求行列式

\[
\left|\begin{array}{llll}
2 & 1 & 1 & 1 \\
1 & 2 & 2 & 1 \\
1 & 2 & 1 & 2 \\
1 & 2 & 2 & 1
\end{array}\right| \text { (另一个版本为 }\left|\begin{array}{llll}
2 & 1 & 1 & 1 \\
1 & 2 & 0 & 1 \\
1 & 0 & 2 & 1 \\
1 & 1 & 1 & 2
\end{array}\right|
\]

的值，并求出第一行元素的代数余子式之和。

\subsection{题目2}
2．（10 分）（可能有误）证明：对任意 \(n\) 阶矩阵 \(A\) ，都有 \(n\) 阶矩阵 \(B, C\) ，使得 \(A=B+C\) ，且 \(B^{\mathrm{T}}=B, C^{\mathrm{T}}=C\) ．

\subsection{题目3}
3．（15 分）在实数域和复数域上将 \(x^{2026}-1\) 分解为不可约因式的乘积。

\subsection{题目4}
4．（15 分）设 \(f(x)=c_{0}+c_{1} x+\cdots+c_{n} x^{n}\) ，证明：若 \(f(x)\) 有 \(n+1\) 个不同的根，那么 \(f(x)\) 是零多项式。

\subsection{题目5}
5．（15 分）将下列向量施密特正交化：\(\alpha_{1}=(1,1)^{\mathrm{T}}, \alpha_{2}=(1,3)^{\mathrm{T}}\) ．

\subsection{题目6}
6．（15 分）（可能有误）设方阵 \(A\) 满足 \(A^{2}-3 A-E=O\) ，证明：\(A+2 E, A-4 E\) 都是可逆矩阵，并用\(A\) 的多项式表示它们的逆．

\subsection{题目7}
7．（20 分）在4维实向量空间中，设

\[
\begin{gathered}
\alpha_{1}=(2,1,0,1), \alpha_{2}=(1,1,1,-1), \alpha_{3}=(3,2,1,0), \\
\beta_{1}=(-1,0,1,2), \beta=(-1,3,7,1) .
\end{gathered}
\]

\(\alpha_{1}, \alpha_{2}, \alpha_{3}\) 生成的子空间为 \(V_{1}, \beta_{1}, \beta_{2}\) 生成的子空间为 \(V_{2}\) 。求 \(V_{1}, V_{2}\) 的维数及 \(V_{1}+V_{2}\) 的一组基。

\subsection{题目8}
8．（25 分）设 \(A\) 是 \(n\) 阶可逆实矩阵，\(E\) 是 \(n\) 阶单位矩阵，证明：正数 \(t\) 足够小时，\(\left(A^{\mathrm{T}} A-t E\right)^{-1}\) 是正定矩阵。

\subsection{题目9}
9．（25 分）已知 \(\cos \frac{\theta}{2} \neq 0\) ，矩阵

\[
U=\left(\begin{array}{ccc}
1 & 0 & 0 \\
0 & \cos \theta & -\sin \theta \\
0 & \sin \theta & \cos \theta
\end{array}\right) .
\]

\(E\) 是 3 阶单位矩阵。证明：\(E+U\) 可逆，且

\[
(E-U)(E+U)^{-1}=\tan \frac{\theta}{2}\left(\begin{array}{ccc}
0 & 0 & 0 \\
0 & 0 & 1 \\
0 & -1 & 0
\end{array}\right) .
\]
\end{document}
